% Options for packages loaded elsewhere
\PassOptionsToPackage{unicode}{hyperref}
\PassOptionsToPackage{hyphens}{url}
\documentclass[
  a4paper,
]{article}
\usepackage{xcolor}
\usepackage{amsmath,amssymb}
\setcounter{secnumdepth}{-\maxdimen} % remove section numbering
\usepackage{iftex}
\ifPDFTeX
  \usepackage[T1]{fontenc}
  \usepackage[utf8]{inputenc}
  \usepackage{textcomp} % provide euro and other symbols
\else % if luatex or xetex
  \usepackage{unicode-math} % this also loads fontspec
  \defaultfontfeatures{Scale=MatchLowercase}
  \defaultfontfeatures[\rmfamily]{Ligatures=TeX,Scale=1}
\fi
\usepackage{lmodern}
\ifPDFTeX\else
  % xetex/luatex font selection
\fi
% Use upquote if available, for straight quotes in verbatim environments
\IfFileExists{upquote.sty}{\usepackage{upquote}}{}
\IfFileExists{microtype.sty}{% use microtype if available
  \usepackage[]{microtype}
  \UseMicrotypeSet[protrusion]{basicmath} % disable protrusion for tt fonts
}{}
\makeatletter
\@ifundefined{KOMAClassName}{% if non-KOMA class
  \IfFileExists{parskip.sty}{%
    \usepackage{parskip}
  }{% else
    \setlength{\parindent}{0pt}
    \setlength{\parskip}{6pt plus 2pt minus 1pt}}
}{% if KOMA class
  \KOMAoptions{parskip=half}}
\makeatother
\setlength{\emergencystretch}{3em} % prevent overfull lines
\providecommand{\tightlist}{%
  \setlength{\itemsep}{0pt}\setlength{\parskip}{0pt}}
\usepackage{bookmark}
\IfFileExists{xurl.sty}{\usepackage{xurl}}{} % add URL line breaks if available
\urlstyle{same}
\hypersetup{
  hidelinks,
  pdfcreator={LaTeX via pandoc}}

\ifXeTeX
\special{pdf:trailerid [ <2bde66321771adeef77e913c4aef2fbc> <2bde66321771adeef77e913c4aef2fbc> ]}
\fi
\ifPDFTeX
\pdftrailerid{}
\pdftrailer{/ID [ <2bde66321771adeef77e913c4aef2fbc> <2bde66321771adeef77e913c4aef2fbc> ]}
\fi
\ifLuaTeX
\pdfvariable trailerid {[ <2bde66321771adeef77e913c4aef2fbc> <2bde66321771adeef77e913c4aef2fbc> ]}
\fi

\author{}
\date{}

\begin{document}

\section{Database Consistency Test}\label{database-consistency-test}

\section{Introduction}\label{introduction}

This EngineeringPaper.xyz sheet shows how to calculated the maximum
stress and displacement for a cantilever beam loaded by a vertical force
at its end. The geometry of the beam is shown in the figure below. The
beam has a rectangular cross section with height h and width b. The
maximum stress in a cantilever with length much larger than its height
is due to bending stress that occurs at the base of the beam. For a
downward force, the maximum stress will be a tensile force at the top of
the beam and a compressive stress at the bottom of the beam. For beam
sections that are symmetric top to bottom, the maximum tensile and
compressive stresses will have equal magnitude. \# \# \# Select Beam
Material

\section{Select Beam Section}\label{select-beam-section}

Source for beam section equations and images:
\url{https://en.wikipedia.org/wiki/List_of_second_moments_of_area}

\section{Define the Other Beam
Parameters}\label{define-the-other-beam-parameters}

\[ l=500\left[mm\right]   \]

\[ F=500\left[N\right]   \]

\section{Define the Deflection and Stress
Equations}\label{define-the-deflection-and-stress-equations}

The maximum displacement of a cantilever beam loaded at its end is
defined by {[}1{]}:

\[ y_{max}=-\frac{F\cdot l^{3}}{3\cdot E\cdot I_{x}}   \]

where Ix is the area moment of inertia of the cross section about the
x-axis as selected in the table above. The maximum displacement can now
be queried:

\[\boxed{ y_{max}=\left[mm\right] =-0.165297794307557 \left[mm\right]  }\]

The maximum bending stress in a beam section is given by:

\[ \sigma_{max}=\frac{M\cdot c}{I_{x}}   \]

where M is the bending moment at the section of interest and c is
distance from the neutral axis of the beam to the top or bottom of the
beam. The maximum moment occurs at the base of the beam and is defined
by:

\[ M=l\cdot F   \]

where l is the length of the beam. Now that everything is defined, the
maximum stress can be queried and converted to {[}MPa{]} units:

\[\boxed{ \sigma_{max}=\left[MPa\right] =5.33333333333333 \left[MPa\right]  }\]

These results can be compared to the results obtained from a finite
element analysis using SolidWorks. The beam has the same dimensions,
loading, and material as the rectangular cross-section and the aluminum
table options above. The displacement plot is shown below. The maximum
displacement magnitude obtained using finite element analysis is .175
{[}mm{]}, which is close to the .165 {[}mm{]} obtained using the
equations above.

Similarly, the maximum stress calculation can be compared to the
SolidWorks results. The maximum stress obtained from SolidWorks is 6.5
{[}MPa{]} as shown below. Note that this is higher than the 5.3
{[}MPa{]} obtained using the beam equation. This is the result of the
stress concentration that occurs at the based of the cantilever where it
meets the fixed boundary condition. A littler further from the fixed
boundary condition, the finite element values are close to the
calculated values.

{[}1{]} Norton, R. L.
``\href{https://www.pearson.com/us/higher-education/program/Norton-Machine-Design-5th-Edition/PGM275676.html}{Machine
design. A integrated approach, 5th Editi.}'' (2013).

\[ y=\frac{F}{6\cdot E\cdot I_{x}}\cdot\left(x^{3}-3\cdot l\cdot x^{2}\right)   \]

\section{All Colors}\label{all-colors}

one

{\textcolor[HTML]{E60000}{two}}

{\textcolor[HTML]{FF9900}{three}}

{\textcolor[HTML]{FFFF00}{four}}

{\textcolor[HTML]{008A00}{five}}

{\textcolor[HTML]{0066CC}{six}}

{\textcolor[HTML]{9933FF}{seven}}

{\textcolor[HTML]{FFFFFF}{eight}}

{\textcolor[HTML]{FACCCC}{nine}}

{\textcolor[HTML]{FFEBCC}{ten}}

{\textcolor[HTML]{FFFFCC}{eleven}}

{\textcolor[HTML]{CCE8CC}{twelve}}

{\textcolor[HTML]{CCE0F5}{thirteen}}

{\textcolor[HTML]{EBD6FF}{fourteen}}

{\textcolor[HTML]{BBBBBB}{fifteen}}

{\textcolor[HTML]{F06666}{sixteen}}

{\textcolor[HTML]{FFC266}{seventeen}}

{\textcolor[HTML]{FFFF66}{eighteen}}

{\textcolor[HTML]{66B966}{nineteen}}

{\textcolor[HTML]{66A3E0}{twenty}}

{\textcolor[HTML]{C285FF}{twenty one}}

{\textcolor[HTML]{888888}{twenty two}}

{\textcolor[HTML]{A10000}{twenty three}}

{\textcolor[HTML]{B26B00}{twenty four}}

{\textcolor[HTML]{B2B200}{twenty five}}

{\textcolor[HTML]{006100}{twenty six}}

{\textcolor[HTML]{0047B2}{twenty seven}}

{\textcolor[HTML]{6B24B2}{twenty eight}}

{\textcolor[HTML]{444444}{twenty nine}}

{\textcolor[HTML]{5C0000}{thirty}}

{\textcolor[HTML]{663D00}{thirty one}}

{\textcolor[HTML]{666600}{thirty two}}

{\textcolor[HTML]{003700}{thirty three}}

{\textcolor[HTML]{002966}{thirty four}}

{\textcolor[HTML]{3D1466}{thirty five}}

{\colorbox[HTML]{FFFF00}{highlight}}

{\textcolor[HTML]{E60000}{\colorbox[HTML]{FFFF00}{color and highlight}}}

\textbf{{\textcolor[HTML]{0066CC}{color and bold}}}

\emph{\textbf{{\textcolor[HTML]{0066CC}{color and bold and italics}}}}

\emph{\textbf{{\textcolor[HTML]{0066CC}{\colorbox[HTML]{FFFFCC}{color
and highlight and bold and italics}}}}}

no text formatting

Created with \href{https://engineeringpaper.xyz}{EngineeringPaper.xyz}

\end{document}
